% ==================================================================
% YangTex · 由 Markdown 模板自动生成（内置 Markdown -> LaTeX 转换器）
% 源模板：md_literature_card.md
% 这是标准 LaTeX：可直接编辑，实时编译 / SyncTeX 双向同步照常可用。
% 需要继续用 Markdown 写时，选中文本后点「MD→TeX 转换」。
% ==================================================================
\documentclass[UTF8,a4paper,11pt,fontset=fandol]{ctexart}
\usepackage{amsmath}
\usepackage{amssymb}
\usepackage{mathtools}
\usepackage{amsthm}
\usepackage{graphicx}
\usepackage{xcolor}
\usepackage{enumitem}
\usepackage{geometry}
\usepackage{hyperref}
\geometry{margin=2.5cm}
\newtheorem{theorem}{定理}[section]
\newtheorem{lemma}[theorem]{引理}
\newtheorem{definition}[theorem]{定义}
\newtheorem{proposition}[theorem]{命题}
\newtheorem{corollary}[theorem]{推论}
\title{文献精读卡}
\author{杨建伟}
\date{\today}
\begin{document}
\maketitle
\section{文献信息}

\begin{itemize}
  \item \textbf{题目}：On the regularity of weak solutions to the Navier--Stokes equations
  \item \textbf{作者 / 年份}：待填写
  \item \textbf{期刊 / 卷页}：待填写
\end{itemize}

\subsection{一句话结论}

本文在 $L^3$ 临界空间中给出了弱解的 \textbf{局部正则性} 判据。

\subsection{主要假设}

\begin{itemize}
  \item 区域 $\Omega \subset \mathbb{R}^3$ 有界且边界光滑
  \item $u_0 \in L^2_\sigma(\Omega)$
  \item 外力 $f \in L^2(0,T;H^{-1})$
\end{itemize}

\subsection{核心估计}

\[
\|u(t)\|_{L^\infty} \le C t^{-1/2} \|u_0\|_{L^3}
\]

\begin{quote}
关键技巧：把非线性项放进 \textit{Carleson 型} 估计，而非直接做能量估计。
\end{quote}

\subsection{我关心的问题}

\begin{enumerate}
  \item 条件中的 $C$ 是否依赖 $\nu$？
  \item 结论能否推广到有界区域的 \textbf{Navier 滑移边界}？
  \item 与 Serrin 条件 $2/p + 3/q \le 1$ 的关系
\end{enumerate}

\subsection{可复用的引理}

\begin{center}
\begin{tabular}{l|l|l}
\hline
引理 & 作用 & 是否需要改写 \\
\hline
Lemma 2.1 & 局部能量不等式 & 需要 \\
Lemma 3.4 & 紧性 & 可直接引用 \\
\hline
\end{tabular}
\end{center}

\subsection{评价}

\begin{itemize}
  \item 证明结构清晰，但第 4 节缩放论证略简
  \item 值得在自己的 PDE 笔记里重写一遍
\end{itemize}

\par\noindent\rule{\textwidth}{0.4pt}\par

精读卡由 YangTex Markdown 模板生成。

\end{document}
